\chapter{वर्गमूल}\label{ch:g9:sqrt}

भुजा $1$ वाले वर्ग का विकर्ण कितना लंबा है? पाइथागोरस प्रमेय जवाब
देता है $\sqrt 2$ --- वह संख्या जिसका वर्ग $2$ है, और जो आख़िर में कोई
\omterm{def:g9:fractions:fraction}{भिन्न} निकलती ही नहीं। यह अध्याय \omterm{def:g9:sqrt:def}{वर्गमूल} की परिभाषा देता है, उनके साथ
गणना के नियम पक्के करता है, और $\sqrt{75}$ जैसे व्यंजकों को सरल करना
सिखाता है।

\section{परिभाषा और पहले गुण}

\begin{definition}[वर्गमूल]\label{def:g9:sqrt:def}
मान लो $a \geq 0$ है। $a$ का \emph{वर्गमूल}\index{वर्गमूल}, जिसे
$\sqrt a$ लिखा जाता है, वह अकेली \emph{अऋणात्मक} संख्या है जिसका वर्ग
$a$ हो:
\[
\sqrt a \geq 0
\qquad\text{और}\qquad
\left(\sqrt a\right)^2 = a .
\]
ऋण संख्याओं का कोई वर्गमूल नहीं होता, क्योंकि हर वर्ग अऋणात्मक होता है।
\end{definition}

\begin{example}
$\sqrt{49} = 7$, $\sqrt 0 = 0$, $\sqrt 1 = 1$,
$\sqrt{2.25} = 1.5$। सबसे पहले जो \omterm{def:g9:sqrt:def}{वर्गमूल} याद कर लेने चाहिए वे
\emph{पूर्ण वर्गों} के हैं: $1, 4, 9, 16, 25, 36, 49, 64, 81, 100,
121, 144$।
\end{example}

\begin{omfigure}
\begin{tikzpicture}[scale=1.35]
  \draw[thick, fill=omDef!8] (0,0) rectangle (2,2);
  \draw[omThm, thick] (0,0) -- (2,2);
  \node[below, font=\small] at (1,0) {$1$};
  \node[left, font=\small] at (0,1) {$1$};
  \node[omThm, above left=-2pt, font=\small, rotate=45] at (1.08,0.92)
    {$\sqrt2$};
  \draw[thin] (1.75,0) -- (1.75,0.25) -- (2,0.25);
\end{tikzpicture}

{\small इकाई वर्ग के विकर्ण की लंबाई पाइथागोरस प्रमेय से
$\sqrt{1^2 + 1^2} = \sqrt2 \approx 1.414$ होती है
(\cref{ch:g9:trig} उसे याद दिलाता है)। यह संख्या अपरिमेय है: इसके
दशमलव कभी नहीं दोहराते।}
\end{omfigure}

\begin{proposition}[वर्ग का वर्गमूल]\label{prop:g9:sqrt:square}
हर वास्तविक संख्या $x$ के लिए (चाहे वह धन हो या न हो):
\[
\sqrt{x^2} = \abs{x} =
\begin{cases}
x & \text{अगर } x \geq 0,\\
-x & \text{अगर } x < 0 .
\end{cases}
\]
\end{proposition}

\begin{proof}
संख्या $\abs{x}$ अऋणात्मक है और उसका वर्ग $x^2$ है (किसी संख्या और
उसकी \omterm{def:g7:negatives:def}{विपरीत संख्या} का वर्ग एक ही होता है)। $x^2$ वर्ग वाली अऋणात्मक
संख्या अकेली होती है, इसलिए वही $\sqrt{x^2}$ है।
\end{proof}

\begin{example}
$\sqrt{(-5)^2} = \sqrt{25} = 5 = \abs{-5}$: \omterm{def:g9:sqrt:def}{वर्गमूल} चिह्न को ``भूल''
जाता है, वापस नहीं लाता। ख़ास तौर पर ऋण $x$ के लिए $\sqrt{x^2} = x$
लिखना \emph{ग़लत} है।
\end{example}

\section{वर्गमूलों के गुणनफल और भागफल}

\begin{theorem}[गुणा और भाग के नियम]\label{thm:g9:sqrt:rules}
सभी $a \geq 0$ और $b \geq 0$ के लिए:
\[
\sqrt{a b} = \sqrt a \times \sqrt b,
\qquad\text{और } b > 0 \text{ के लिए:}\quad
\sqrt{\frac ab} = \frac{\sqrt a}{\sqrt b} .
\]
\end{theorem}

\begin{proof}
संख्या $\sqrt a \times \sqrt b$ अऋणात्मक है (अऋणात्मक संख्याओं का
\omterm{def:g2:mult:def}{गुणनफल}), और उसका वर्ग है
\[
\left(\sqrt a \times \sqrt b\right)^2
= \left(\sqrt a\right)^2 \times \left(\sqrt b\right)^2
= a b .
\]
$ab$ वर्ग वाली अऋणात्मक संख्या अकेली होती है, इसलिए वह $\sqrt{ab}$ के
बराबर है। भाग वाला नियम भी इसी तरह सिद्ध होता है।
\end{proof}

\begin{remark}[योग के लिए ऐसा कोई नियम नहीं!]
$\sqrt{a + b}$ आम तौर पर $\sqrt a + \sqrt b$ \emph{नहीं} होता:
\[
\sqrt{9 + 16} = \sqrt{25} = 5,
\qquad\text{पर}\qquad
\sqrt 9 + \sqrt{16} = 3 + 4 = 7 .
\]
\end{remark}

\begin{method}[$\sqrt n$ को सरल करना]\label{met:g9:sqrt:simplify}
किसी पूर्णांक के \omterm{def:g9:sqrt:def}{वर्गमूल} को सरल करने के लिए:
\begin{enumerate}
  \item $n$ का \omterm{def:g9:arith:divisor}{भाजक} बनने वाला सबसे बड़ा पूर्ण वर्ग ढूँढ़ो (ज़रूरत हो तो
        $n$ का गुणनखंडन कर लो);
  \item \omterm{def:g2:mult:def}{गुणनफल} वाले नियम से टुकड़े करो और वर्ग बाहर निकालो:
        $\sqrt{k^2 m} = k\sqrt m$;
  \item जाँचो कि मूल के नीचे जो बचा है उसका कोई पूर्ण वर्ग \omterm{def:g9:arith:divisor}{भाजक} न
        रह गया हो।
\end{enumerate}
\end{method}

\begin{example}\label{ex:g9:sqrt:simplify}
$\sqrt{75}$ को सरल करो: चूँकि $75 = 25 \times 3$ है,
\[
\sqrt{75} = \sqrt{25 \times 3} = \sqrt{25} \times \sqrt 3 = 5\sqrt3 .
\]
$\sqrt{72}$ को सरल करो: चूँकि $72 = 36 \times 2$ है, इसलिए
$\sqrt{72} = 6\sqrt2$। और एक \omterm{def:g3:division:remainder}{भागफल}:
$\sqrt{\dfrac{49}{16}} = \dfrac{\sqrt{49}}{\sqrt{16}} = \dfrac74$।
\end{example}

\begin{example}[वर्गमूल जोड़ना]\label{ex:g9:sqrt:add}
वर्गमूलों के \omterm{def:g1:addition:def}{योगफल} तभी सरल होते हैं जब मूल \emph{एक जैसे} हों।
$\sqrt{12} + \sqrt{27} - \sqrt{48}$ निकालो, और पहले हर पद सरल करो:
\begin{align*}
\sqrt{12} &= \sqrt{4 \times 3} = 2\sqrt3, \\
\sqrt{27} &= \sqrt{9 \times 3} = 3\sqrt3, \\
\sqrt{48} &= \sqrt{16 \times 3} = 4\sqrt3,
\end{align*}
इसलिए \omterm{def:g1:addition:def}{योगफल} $2\sqrt3 + 3\sqrt3 - 4\sqrt3 = (2 + 3 - 4)\sqrt3 = \sqrt3$ है।
\end{example}

\begin{example}[वर्गमूलों के साथ प्रसार]\label{ex:g9:sqrt:expand}
बीजगणित की सर्वसमिकाएँ वर्गमूलों पर भी लगती हैं।
$(a+b)^2 = a^2 + 2ab + b^2$ से $\left(\sqrt5 + 2\right)^2$ का प्रसार करो:
\[
\left(\sqrt5 + 2\right)^2 = 5 + 2 \times 2\sqrt5 + 4 = 9 + 4\sqrt5 .
\]
और तीसरी सर्वसमिका $(a+b)(a-b) = a^2 - b^2$ से:
\[
\left(\sqrt7 + \sqrt3\right)\left(\sqrt7 - \sqrt3\right) = 7 - 3 = 4 :
\]
दो अपरिमेय संख्याओं का \omterm{def:g2:mult:def}{गुणनफल} पूर्णांक हो सकता है।
\end{example}

\section{समीकरण \texorpdfstring{$x^2 = a$}{x2 = a}}

\begin{theorem}[$x^2 = a$ हल करना]\label{thm:g9:sqrt:equation}
\begin{enumerate}
  \item अगर $a > 0$ हो, तो \omterm{def:g8:equations:def}{समीकरण} $x^2 = a$ के ठीक दो हल होते हैं:
        $\sqrt a$ और $-\sqrt a$;
  \item अगर $a = 0$ हो, तो अकेला हल $0$ है;
  \item अगर $a < 0$ हो, तो कोई हल नहीं है।
\end{enumerate}
\end{theorem}

\begin{proof}
$a \geq 0$ होने पर $x^2 = a$ को $x^2 - \left(\sqrt a\right)^2 = 0$
लिखो, और वर्गों के \omterm{ex:g1:subtraction:difference}{अंतर} की तरह गुणनखंड करो:
\[
\left(x - \sqrt a\right)\left(x + \sqrt a\right) = 0 ,
\]
इसलिए $x = \sqrt a$ या $x = -\sqrt a$ ($a = 0$ पर ये दोनों एक ही हो
जाते हैं)। और $a < 0$ पर कोई हल नहीं है, क्योंकि हर $x$ के लिए
$x^2 \geq 0 > a$ होता है।
\end{proof}

\begin{omfigure}
\begin{tikzpicture}
\begin{axis}[omaxis,
  width=7.6cm, height=5.4cm,
  xmin=-3.1, xmax=3.1, ymin=-1.6, ymax=5.4,
  xtick={-1.732,1.732}, xticklabels={$-\sqrt3$,$\sqrt3$},
  ytick={3}, yticklabels={$3$}]
  \addplot[omDef, thick, domain=-2.25:2.25, samples=90] {x^2};
  \addplot[omThm, thick, domain=-2.9:2.9] {3};
  \fill (axis cs:-1.732,3) circle (1.5pt);
  \fill (axis cs:1.732,3) circle (1.5pt);
  \draw[dashed] (axis cs:-1.732,3) -- (axis cs:-1.732,0);
  \draw[dashed] (axis cs:1.732,3) -- (axis cs:1.732,0);
\end{axis}
\end{tikzpicture}

{\small परवलय $y = x^2$ पर $x^2 = 3$ हल करना: आड़ी रेखा $y = 3$ उसे
$x = -\sqrt3$ और $x = \sqrt3$ पर काटती है।}
\end{omfigure}

\begin{example}
$x^2 = 16$: हल $4$ और $-4$। \quad
$x^2 = 5$: हल $\sqrt5$ और $-\sqrt5$ (ठीक-ठीक मान; $\sqrt5
\approx 2.236$)। \quad
$x^2 = -9$: कोई हल नहीं। \quad
$3x^2 = 21$: पहले $3$ से भाग दो, $x^2 = 7$, हल $\pm\sqrt7$।
\end{example}

\section{अभ्यास}

\begin{exercise}[$\star$]\label{exo:g9:sqrt:1}
कैलकुलेटर के बिना निकालो:
\[
\sqrt{64}, \qquad \sqrt{100}, \qquad \sqrt{0.09}, \qquad
\sqrt{\tfrac{1}{25}}, \qquad \left(\sqrt{13}\right)^2, \qquad
\sqrt{(-4)^2} .
\]
\end{exercise}

\begin{exercise}[$\star$]\label{exo:g9:sqrt:2}
सरल करो:
\[
\sqrt{50}, \qquad \sqrt{45}, \qquad \sqrt{98}, \qquad \sqrt{300} .
\]
\end{exercise}

\begin{exercise}[$\star$]\label{exo:g9:sqrt:3}
निकालो और सरल करो:
\[
\sqrt2 \times \sqrt{18}, \qquad
\frac{\sqrt{75}}{\sqrt3}, \qquad
\sqrt5 \times \sqrt{20}, \qquad
\frac{\sqrt{8}}{\sqrt{50}} .
\]
\end{exercise}

\begin{exercise}[$\star$]\label{exo:g9:sqrt:4}
हल करो: $x^2 = 36$; \quad $x^2 = 11$; \quad $x^2 + 4 = 0$; \quad
$2x^2 = 50$।
\end{exercise}

\begin{exercise}[$\star$]\label{exo:g9:sqrt:5}
एक ही पद में समेटो: $\sqrt{20} + \sqrt{45}$, फिर
$3\sqrt8 - \sqrt{18} + \sqrt2$।
\end{exercise}

\begin{exercise}[$\star\star$]\label{exo:g9:sqrt:6}
प्रसार करो और सरल करो:
\[
\left(\sqrt3 + 1\right)^2, \qquad
\left(2\sqrt5 - 3\right)^2, \qquad
\left(\sqrt6 + \sqrt2\right)\left(\sqrt6 - \sqrt2\right).
\]
\end{exercise}

\begin{exercise}[$\star\star$]\label{exo:g9:sqrt:7}
एक वर्गाकार खेत का \omterm{def:g6:measure:area}{क्षेत्रफल} $6400$ m$^2$ है। उसकी भुजा कितनी लंबी
है? और उसका विकर्ण (पहले ठीक-ठीक मान, फिर निकटतम मीटर तक गोल किया
हुआ)?
\end{exercise}

\begin{exercise}[$\star\star$]\label{exo:g9:sqrt:8}
दिखाओ कि $\dfrac{1}{\sqrt2} = \dfrac{\sqrt2}{2}$ है (\omterm{def:g4:fractions:def}{अंश} और हर दोनों
को $\sqrt2$ से गुणा करो)। इसी तरकीब से $\dfrac{6}{\sqrt3}$ और
$\dfrac{10}{\sqrt5}$ को हर में \omterm{def:g9:sqrt:def}{वर्गमूल} के बिना लिखो।
\end{exercise}

\begin{exercise}[$\star\star$]\label{exo:g9:sqrt:9}
सही या ग़लत? प्रमाण या प्रतिउदाहरण से कारण बताओ।
\begin{enumerate}
  \item सभी $a, b \geq 0$ के लिए: $\sqrt{ab} = \sqrt a \sqrt b$।
  \item सभी $a, b \geq 0$ के लिए: $\sqrt{a+b} = \sqrt a + \sqrt b$।
  \item सभी $x$ के लिए: $\sqrt{x^2} = x$।
  \item $\left(3\sqrt2\right)^2 = 18$।
\end{enumerate}
\end{exercise}

\begin{exercise}[$\star\star\star$]\label{exo:g9:sqrt:10}
मान लो $x = \sqrt{7 + 4\sqrt3}$ है। $\left(2 + \sqrt3\right)^2$
निकालो, और इससे $x$ का सरल रूप निकालो। वही सवाल
$\sqrt{9 - 4\sqrt5}$ के लिए ($\left(\sqrt5 - 2\right)^2$ जैसे वर्ग
की ओर बढ़ो, और चिह्न का ध्यान रखो)।
\end{exercise}

\section{समस्या: वह संख्या जो भिन्न नहीं है}

\begin{problem}[{सप्ताहांत समस्या --- $\sqrt2$ के अपरिमेय होने का वह
पौराणिक प्रमाण, उसके कई भाई-बंधु, और उसका लगभग मान निकालने की हेरॉन
की हैरान कर देने वाली अच्छी तरकीब}]
\label{pb:g9:sqrt:1}
इकाई वर्ग का विकर्ण एकदम असली लंबाई है --- तुमने उसे
\cref{ex:g8:pythagoras:diagonal} में खींचा भी था --- और फिर भी, जैसा
पाइथागोरस के अनुयायियों ने क़रीब पच्चीस सदी पहले अपने आतंक के साथ
खोजा, \emph{कोई भी \omterm{def:g9:fractions:fraction}{भिन्न}} उसे नाप नहीं सकती। किंवदंती कहती है कि इस
खोज की सज़ा डुबोकर दी गई थी। यह समस्या तुम्हें उस अमर प्रमाण से गुज़ार
देती है (चार सवाल, और वह ज़िंदगी भर के लिए तुम्हारा), शिकारों की
\omterm{def:g7:stats:frequency}{गिनती} बढ़ाती है, और उस तरकीब पर ख़त्म होती है जिससे कारीगर दो हज़ार साल
तक इस बेक़ाबू संख्या को क़ाबू में करते रहे।

\medskip
\noindent\textbf{भाग I --- वह प्रमाण।}
\begin{enumerate}
  \item पहले उसे घेरो: सीमाओं के वर्ग निकालकर जाँचो कि
        $1.4 < \sqrt2 < 1.5$ है, फिर कि $1.41 < \sqrt2 < 1.42$ है।
        एक अंक और: $\sqrt2$ तीन दशमलव वाली किन दो संख्याओं के बीच
        है?
  \item अब --- विरोधाभास के लिए --- मान लो कि किसी \emph{सबसे सरल
        रूप} वाली \omterm{def:g9:fractions:fraction}{भिन्न} (\cref{met:g9:arith:simplify}) के लिए
        $\sqrt2 = \frac ab$ है। दोनों पक्षों का वर्ग करो और दिखाओ
        कि $a^2 = 2 b^2$ है। $a^2$ सम है या विषम?
  \item \cref{pb:g8:pythagoras:1} के सम-विषम तथ्य कहते हैं: \omterm{def:g3:numbers:evenodd}{विषम
        संख्याओं} के वर्ग विषम होते हैं। इससे नतीजा निकालो कि $a$ सम
        है, $a = 2k$ लिखो, रखकर देखो --- और दिखाओ कि $b$ को भी सम
        होना पड़ेगा।
  \item विरोधाभास कहाँ है? नतीजा निकालो, और प्रमेय पूरी तरह लिखो:
        \emph{$\sqrt2$ अपरिमेय है --- वह पूर्ण संख्याओं का कोई
        \omterm{def:g3:division:remainder}{भागफल} नहीं है}।
  \item प्रमाण एक बार, चुपके से, ``सबसे सरल रूप'' शब्दों पर टिका
        था। ठीक वह क़दम बताओ जो उनके बिना ढह जाता, और यह भी समझाओ
        कि सबसे सरल रूप मान लेना पहले से ही जायज़ क्यों था।
\end{enumerate}

\medskip
\noindent\textbf{भाग II --- शिकार बढ़ते जाते हैं।}
\begin{enumerate}[resume]
  \item अभाज्य गुणनखंडनों वाली दूसरी शैली का प्रमाण
        (\cref{thm:g9:arith:factorization}): $a^2 = 3b^2$ में दोनों
        पक्षों पर \emph{$3$ के \omterm{def:g8:powers:def}{घातांक}} का सम-विषम होना भिड़ाओ
        (\cref{pb:g9:arith:1}: वर्गों के \omterm{def:g8:powers:def}{घातांक} सम होते हैं)।
        नतीजा निकालो कि $\sqrt3$ अपरिमेय है।
  \item किसी आम पूर्ण संख्या $n$ के लिए $a^2 = n b^2$ पर वही \omterm{def:g8:powers:def}{घातांक}
        वाला तर्क चलाओ: किन $n$ के लिए वह विरोधाभास खड़ा करता है,
        और किनके लिए नाकाम रहता है? पूरा नतीजा लिखो: $\sqrt n$ ठीक
        तब अपरिमेय होता है जब\,\dots
  \item ढाल वाली प्रमेयिका सिद्ध करो: किसी शून्येतर परिमेय संख्या
        और किसी अपरिमेय संख्या का \omterm{def:g2:mult:def}{गुणनफल} अपरिमेय होता है। (मान लो
        $r x = q$ है, जहाँ $r, q$ परिमेय हैं और $r \neq 0$; अब $x$
        निकालो।) इससे नतीजा निकालो कि $2\sqrt2$ और
        $\frac{\sqrt2}{2}$ अपरिमेय हैं।
  \item सिद्ध करो कि $1 + \sqrt2$ अपरिमेय है। फिर वह सुंदर वाला:
        मान लो $s = \sqrt2 + \sqrt3$ परिमेय होता; $s^2$ निकालो,
        $\sqrt6$ को अलग करो, और विरोधाभास ढूँढ़ो।
  \item अब जोश थोड़ा ठंडा करो: दो ऐसी अपरिमेय संख्याएँ बताओ जिनका
        \emph{\omterm{def:g1:addition:def}{योगफल}} परिमेय हो, और दो ऐसी जिनका \emph{\omterm{def:g2:mult:def}{गुणनफल}}
        परिमेय हो। (अंकगणित अपरिमेयता को बचाकर नहीं रखता --- हर
        हालत का अपना प्रमाण चाहिए।)
\end{enumerate}

\medskip
\noindent\textbf{भाग III --- हेरॉन की तरकीब।}
कैलकुलेटरों से दो हज़ार साल पहले सिकंदरिया के हेरॉन ने $\sqrt2$ का
लगभग मान ऐसे निकाला: \emph{कोई $x$ अंदाज़ो; उस अंदाज़े की जगह $x$ और
$\frac2x$ का औसत रख दो; यही दोहराते रहो।}
\begin{enumerate}[resume]
  \item अंदाज़े $x = 1$ से शुरू करो और अगले दो अंदाज़े ठीक-ठीक भिन्नों
        के रूप में निकालो।
  \item तीसरा अंदाज़ा भी ठीक-ठीक \omterm{def:g9:fractions:fraction}{भिन्न} के रूप में निकालो, और उसका
        दशमलव मान भी। सवाल 1 से तुलना करो: $\sqrt2$ के कितने दशमलव
        अभी से सही हैं?
  \item तरकीब के पीछे का विचार: \omterm{def:g6:measure:area}{क्षेत्रफल} $2$ वाले जिस आयत की एक
        भुजा $x$ है, उसकी दूसरी भुजा $\frac2x$ होती है। दिखाओ कि
        $\sqrt2$ हमेशा $x$ और $\frac2x$ के \emph{बीच} रहता है
        (दोनों हालतें $x^2 > 2$ और $x^2 < 2$ देखो), इसलिए दोनों
        भुजाओं का औसत लेना आयत को वर्ग की ओर दबाता है।
  \item अंदाज़ों $\frac32$, $\frac{17}{12}$, $\frac{577}{408}$ में
        एक रत्न छिपा है: $3^2 - 2 \times 2^2$ निकालो, फिर
        $17^2 - 2 \times 12^2$, फिर $577^2 - 2 \times 408^2$।
        भाग I ने सिद्ध किया कि $a^2 - 2b^2 = 0$ नामुमकिन है ---
        हेरॉन की भिन्नें उस नामुमकिन के कितने पास पहुँचती हैं?
  \item समापन, दो वाक्यों में: इकाई वर्ग का विकर्ण काग़ज़ पर मौजूद
        है, कोई \omterm{def:g9:fractions:fraction}{भिन्न} उसे नापती नहीं, और उसका दशमलव लेखन कभी दोहरा
        नहीं सकता (\cref{pb:g9:fractions:1})। तो संख्या रेखा में
        कैसी ``नई संख्याएँ'' होनी ही चाहिए --- और इस श्रृंखला में
        उनकी पूरी कहानी कहाँ कही गई है?

\end{enumerate}
\end{problem}
